\section{Background}
\label{sec:background}

% CONTENT GUIDANCE (delete before submission):
% - Subsection 2.1: KV cache in Transformer serving. PagedAttention is
%   the obvious citation (\cite{kwon2023pagedattention}); vLLM-style
%   block-based allocation. Optional: FlashAttention's IO-awareness
%   framing (\cite{dao2023flashattention2}, \cite{shah2024flashattention3}).
% - Subsection 2.2: SSM state caching. The d_state alignment and
%   per-layer constant-size cache shape, contrasted with the KV cache's
%   per-token shape. Cite Mamba (\cite{gu2023mamba}, \cite{dao2024mamba2}).
% - Subsection 2.3: hybrid architectures. Jamba (\cite{lieber2024jamba}),
%   Zamba2 (\cite{glorioso2024zamba2}), Samba (\cite{ren2024samba}),
%   Hymba (\cite{dong2024hymba}), RecurrentGemma
%   (\cite{botev2024recurrentgemma}). Note that none of them prescribe a
%   serving-time memory-layout strategy.
% - Source: RFC 0001 sections 2 and 3.1.

\subsection{Hybrid Mamba-Transformer Architectures}
\label{sec:background:hybrid}

Standard transformer attention requires $O(n^2)$ compute for a sequence of length $n$, and the KV cache grows linearly with $n$ during decoding. Conversely, State Space Models (SSMs) like Mamba operate with $O(n)$ compute and require an $O(1)$ cache per layer \cite{gu2023mamba, dao2024mamba2}. The SSM state size remains fixed and strictly independent of the sequence length. Hybrid architectures combine these two approaches to combine SSM efficiency with attention reasoning.

For example, Jamba mixes attention and Mamba layers at a 1:7 ratio, supporting a 256K token context window with approximately 12B active parameters in a Mixture-of-Experts (MoE) configuration \cite{lieber2024jamba}. Other variations adopt similar hybrid structures with different attention-to-SSM ratios \cite{glorioso2024zamba2, ren2024samba, dong2024hymba, botev2024recurrentgemma}. This architectural combination forces inference engines to manage two fundamentally different memory cache types within a single model.

\subsection{KV Cache and SSM State Management}
\label{sec:background:cache}

Standard transformer inference manages the Key-Value (KV) cache using paged memory allocation \cite{kwon2023pagedattention}, often combined with IO-aware attention kernels \cite{dao2022flashattention, shah2024flashattention3}. The KV cache requires a system that handles variable block sizes, supports append-only operations during decoding, and allows dynamic eviction or recomputation when memory pressure increases. In contrast, the SSM state maintains a fixed size per layer, strictly defined by \texttt{state\_dim} multiplied by \texttt{bytes\_per\_element}. This state updates in place and persists across tokens.

Unlike KV blocks, the SSM state cannot be paged because it lacks a temporal token structure to evict \cite{dao2024mamba2}. It requires a single, contiguous allocation per layer per sequence. Consequently, the memory access patterns diverge significantly: KV cache operations rely on paged scatter-gather memory lookups, while SSM state operations require contiguous read and write access. This structural asymmetry is the root cause of fragmentation and overestimation in existing cache pool designs.

\subsection{Limits of Existing Pool Designs}
\label{sec:background:pool}

Current inference systems attempt to manage hybrid models using either unified or static dual-pool designs, both of which fail in two ways. Unified pools, such as the vLLM \texttt{HybridKVCacheCoordinator} approach \cite{kwon2023pagedattention}, force the SSM state to pad up to the attention page size. As reported in vLLM issue \#37121, this padding causes a 7.3$\times$ KV cache capacity overestimation on Qwen3.5-4B-AWQ, leaving 13.7\% effective VRAM utilization at peak memory. The root cause is the system capacity estimator, which incorrectly multiplies the $O(1)$ SSM state by the sequence token count.

Alternatively, static dual-pool architectures like SGLang pre-allocate two physical regions (\texttt{HybridReqToTokenPool} and \texttt{HybridLinearKVPool}) at engine startup \cite{zheng2024sglang}. This design assigns memory using a fixed parameter, such as a \texttt{mamba\_full\_memory\_ratio} defaulting to 0.9. The system cannot rebalance capacity between pools without a full process restart. When the prompt distribution shifts, static pools perform poorly. As we demonstrate in \S\ref{sec:eval}, a default 0.9 ratio triggers 1221 OOM events compared to 552 OOMs at a 0.5 ratio on identical workloads. Neither unified padding nor static allocation handles runtime workload shifts dynamically, motivating a memory management system that simultaneously supports asymmetric cache types and adapts to shifting prompt distributions.
